\chapter*{Vorwort}
\addcontentsline{toc}{chapter}{Vorwort}

\begin{merksatz}
\boxtitel{Die Frage, die Band 22 offenließ}
Das letzte Kapitel von Band 22 endete mit einer Feststellung: Ein Verfahren, das gut vorhersagt, hat nicht verstanden, warum. Die Merkmale, die es benutzt, beschreiben die Daten --- nicht die Krankheit.

\begin{center}
\begin{minipage}{0.9\textwidth}
\textbf{Dieser Band behandelt die andere Frage: Was geschieht, wenn man eingreift?}
\end{minipage}
\end{center}

\begin{center}
\begin{minipage}{0.88\textwidth}
\kursiv{Sie ist die eigentliche Frage der Medizin. \textbf{Kein Arzt will wissen, welche Patienten sterben werden --- er will wissen, was er tun soll}, damit sie es nicht tun. Das ist keine Frage nach Zusammenhängen, sondern nach Wirkungen, und bis vor wenigen Jahrzehnten galt sie als außerhalb der Statistik liegend.}
\end{minipage}
\end{center}
\end{merksatz}

\begin{merksatz}
\boxtitel{Wo der Band steht}
\begin{center}
{\small
\renewcommand{\arraystretch}{1.52}
\begin{tabular}{@{}P{2.2cm}P{6.4cm}@{}}
\toprule
\textbf{Band 17} & Statistik schätzt Größen einer Verteilung und prüft Hypothesen darüber.\\
\textbf{Band 22} & Das Lernen sagt vorher, was bei einem neuen Fall zu erwarten ist.\\
\textbf{Band 21} & Die Spieltheorie fragt, was geschieht, wenn andere auf eine Änderung reagieren.\\
\textbf{Hier} & Die Frage lautet: Was wäre geschehen, wenn man anders gehandelt hätte?\\
\bottomrule
\end{tabular}}
\end{center}

\begin{center}
\begin{minipage}{0.88\textwidth}
\kursiv{\textbf{Die letzte Zeile enthält eine Größe, die niemand beobachten kann.} Ein Patient wurde früh behandelt oder spät, nicht beides. Was mit ihm im anderen Fall geschehen wäre, ist keine fehlende Messung, sondern eine, die es nicht geben kann. Dass sich über solche Größen dennoch rechnen lässt --- unter angebbaren Annahmen ---, ist der Inhalt dieses Bandes.}
\end{minipage}
\end{center}
\end{merksatz}

\begin{merksatz}
\boxtitel{Drei Ergebnisse, die man nicht erwartet}
\begin{center}
{\small
\renewcommand{\arraystretch}{1.55}
\begin{tabular}{@{}P{1.0cm}P{7.4cm}@{}}
\toprule
$1$ & Dieselbe Tabelle von Zahlen ist mit \kursiv{entgegengesetzten} Wirkungen vereinbar. Keine noch so große Datenmenge kann zwischen ihnen entscheiden --- nur Annahmen können es. \kursiv{(Kapitel 1)}\\
$2$ & Für manche Merkmale ist es ein Fehler, sie herauszurechnen. Wer nach allem anpasst, was er gemessen hat, erzeugt Verzerrungen, die vorher nicht da waren. \kursiv{(Kapitel 5)}\\
$3$ & Aus reinen Beobachtungsdaten lässt sich unter prüfbaren Bedingungen dasselbe berechnen wie aus einer randomisierten Studie --- und die Bedingungen sind erstaunlich oft erfüllt. \kursiv{(Kapitel 6 und 11)}\\
\bottomrule
\end{tabular}}
\end{center}

\begin{center}
\begin{minipage}{0.88\textwidth}
\kursiv{Das zweite Ergebnis widerspricht der verbreiteten Übung, in ein Modell möglichst viele Störgrößen aufzunehmen. \textbf{Es gibt Merkmale, deren Berücksichtigung schadet} --- und man erkennt sie nicht an den Daten, sondern nur an einer Überlegung darüber, wie die Daten entstanden sind.}
\end{minipage}
\end{center}
\end{merksatz}

\begin{merksatz}
\boxtitel{Die sechs Teile}
\begin{center}
{\small
\renewcommand{\arraystretch}{1.5}
\begin{tabular}{@{}P{1.4cm}P{3.4cm}P{3.2cm}@{}}
\toprule
\textbf{Teil} & \textbf{Thema} & \textbf{Kapitel}\\
\midrule
I & Das Problem & $1$--$3$\\
II & Graphen und Regeln & $4$--$6$\\
III & Anpassen & $7$--$8$\\
IV & Wenn Anpassen nicht reicht & $9$--$11$\\
V & Über die Zeit, im Einzelfall & $12$--$14$\\
VI & Stand & $15$\\
\bottomrule
\end{tabular}}
\end{center}

\schritt{Der Leitfall}
Am Klinikum Neustadt zeigen die Behandlungsdaten: \kursiv{Patienten, die früh Antibiotika erhielten, starben häufiger.}

\begin{center}
\begin{minipage}{0.88\textwidth}
\kursiv{Die Leitlinien verlangen eine möglichst frühe Gabe, und die Evidenz dafür ist gut. \textbf{Die eigenen Daten des Hauses scheinen das Gegenteil zu zeigen} --- und die Auflösung dieses Widerspruchs führt durch den ganzen Band. Sie beginnt in Kapitel~1 mit einer Tabelle und endet in Kapitel~11 mit einem Verfahren, das aus Routinedaten eine Studie nachbildet.}
\end{minipage}
\end{center}
\end{merksatz}

\begin{warnung}
\boxtitelr{Drei Warnungen}
\begin{center}
{\small
\renewcommand{\arraystretch}{1.55}
\begin{tabular}{@{}P{2.2cm}P{6.4cm}@{}}
\toprule
\textbf{Erste} & Jede kausale Aussage beruht auf Annahmen, die sich aus den Daten nicht prüfen lassen. Das Verfahren macht sie nicht überflüssig --- es macht sie sichtbar.\\
\textbf{Zweite} & Ein Durchschnittseffekt sagt nichts über den einzelnen Patienten. Er kann positiv sein, während er bei der Hälfte schadet.\\
\textbf{Dritte} & Was in einer Bevölkerung wirkt, muss in einer anderen nicht wirken --- und was unter Studienbedingungen wirkt, nicht im Alltag.\\
\bottomrule
\end{tabular}}
\end{center}

\begin{center}
\begin{minipage}{0.88\textwidth}
\kursiv{Die erste Warnung ist die wichtigste und wird oft für einen Einwand gehalten. \textbf{Sie ist das Gegenteil}: Die Stärke der Verfahren dieses Bandes liegt darin, dass sie die Annahmen benennen, statt sie zu verstecken. Wer eine Regression rechnet und von \enquote{Einfluss} spricht, hat dieselben Annahmen gemacht --- nur ohne sie aufzuschreiben.}
\end{minipage}
\end{center}
\end{warnung}

\begin{merksatz}
\boxtitel{Was dieser Band voraussetzt}
\begin{center}
{\small
\renewcommand{\arraystretch}{1.48}
\begin{tabular}{@{}P{3.0cm}P{1.6cm}P{3.0cm}@{}}
\toprule
\textbf{Werkzeug} & \textbf{Band} & \textbf{gebraucht in}\\
\midrule
Bedingte Wahrscheinlichkeit & $4$ & durchgehend\\
Schätzen, Konfidenz & $17$ & Kapitel 3, 7, 8\\
Regression & $17$, $22$ & Kapitel 7, 9, 14\\
Graphen, Wege & $16$ & Kapitel 4--6\\
Lernverfahren & $22$ & Kapitel 8, 14\\
\bottomrule
\end{tabular}}
\end{center}

\begin{center}
\begin{minipage}{0.88\textwidth}
\kursiv{\textbf{Teil~I und II kommen mit Band 0 und den Grundzügen der Wahrscheinlichkeitsrechnung aus.} Die Graphen aus Band 16 werden gebraucht, aber nur in ihren einfachsten Eigenschaften --- gerichtete Kanten, Wege, Kreisfreiheit. Wo ein Satz aus einem Vorgängerband benutzt wird, steht die Fundstelle dabei.}
\end{minipage}
\end{center}
\end{merksatz}
